\section{Fully paid padding and the limits of the motion test}\label{sec:padding}

A native-machine prefix of linear cost enforces the original cut inequality
through any payload computation. For uniformly linear-time payload families,
the resulting runs also satisfy the full limiting conclusions. This supplies
the required motion for the construction program and identifies what those
conditions alone cannot exclude.

\subsection{The exact geometric contract}

We recall the discrete condition from~\cite{FirstPaper2026}. Let $k$ unit-step
work-head paths $h_p:\N\to\Z$ run for $T$ steps. At dimension $d$, take
$r=\ceil{db(d)/G}$ and split time into $N=\ceil{T/r}$ slots. A block is a
length-$r$ interval of tape cells. The original \emph{pair residency} of a slot
is the union, over every step of that slot before $T$, of the current block
and its previous block under the original last-change convention. A chord
joins resident slots $x<z$ only when $x+1<z$ and the block is absent from
every intervening slot: it connects the last slot of one residency visit to
the first slot of the next separated visit. Adjacent resident slots do not
form chords. The theorem uses this original residency, not an arbitrary
collection of convenient edges.

Give the slots labels by $\lambda:\N\to\N$, with at most $m$ slots of any
one label among $0,\ldots,N-1$. Let $B_\lambda$ be the number of adjacent label changes and
$\chi_p(\lambda)$ the number of residency chords of head $p$ whose endpoints
have different labels. The \emph{zero-slack cut inequality} is
\begin{equation}\label{eq:zero-cut}
 (d-B_\lambda)_+
 (d-2m)_+b(d)
 \le rG\left(3kB_\lambda+\sum_{p<k}\chi_p(\lambda)\right)
\end{equation}
for every such labelling and $m$, where $x_+=\max(x,0)$.
It is precisely \texttt{CutIneq d G r T 0 0 h}; the two zero parameters are
the entropy-deficit and boundary-information slack, not a restriction on the
payload's input distribution.
\upformal{Transpose.Tape.Defs}{Transpose.Tape.pairRes,Transpose.Tape.labelChanges,Transpose.Tape.crossChords}
\upformal{Transpose.Tape.DefsM}{Transpose.Tape.CutIneq}

For a linear time bound $T_i\le CP(d_i)$, freeze each path at $T_i$ and put
\[
 p_{i,t}(u)=\frac{1}{d_i}
  \left\lfloor\frac{h_{i,t}(\min(\floor{ud_i}r_i,T_i))}{r_i}\right\rfloor,
 \qquad 0\le u\le c_0:=CG+1.
\]
These are the original block trajectories. The full limiting conclusion means
that some subsequence converges uniformly, and \emph{every} uniformly
convergent subsequence has $1$-Lipschitz limit paths satisfying the original
$M^*$ inequality for every $\theta\in[0,\infty]$, every finite union $S$ of
closed intervals, and every piecewise-constant-off-$S$ labelling whose
off-$S$ classes each have measure at most $\theta$. The inequality is
\eqref{eq:motion-star} with this full range of $\theta$. In addition,
their total integrated level excess is at least one, the sum of their range
measures plus one is at most $kc_0$, and at least one limit path is neither
monotone nor antitone.
\formal{NativeRunLimits}{TransposeFollowup.NativeRunLimits.Conclusion,TransposeFollowup.NativeRunLimits.of_cut}

\subsection{The wrapper theorem}

\begin{theorem}[Universal fully paid padding]\label{thm:padding}
There are absolute integers $A,H,K_0,E\ge0$ such that for every native
machine $M$ there is a single native machine $M'$ with
\begin{align*}
 \#\mathrm{states}(M')&=H+\#\mathrm{states}(M),\\
 k(M')&=K_0+(k(M)+E)\max(1,\operatorname{size}(g(M))),\\
 \operatorname{cellBits}(M')&=3.
\end{align*}
Here $g(M)+1$ is the work-alphabet size and $\operatorname{size}$ is the
binary length of a natural number, with $\operatorname{size}(0)=0$.
The following hold for this same $M'$.
\begin{enumerate}[label=(\alph*)]
\item For every $d\ge1$, every word $w$ of length $P(d)$, every output
$v$ and every $t$ for which $M$ computes $v$ on $w$ within $t$ steps, there
is an actual halting time $T\le AP(d)+t$ of $M'$ on $w$, with output exactly
$v$, whose entire execution satisfies \eqref{eq:zero-cut} with $G=3$ and
$k=k(M')$.
\item For every sequence $d_i\to\infty$ of positive dimensions, words
$w_i$ of length $P(d_i)$, outputs $v_i$, and $a\in\N$ such that $M$ computes
$v_i$ on $w_i$ within $aP(d_i)$ steps, there are actual halting times
$T_i\le(A+a)P(d_i)$ for $M'$ with those same outputs and the cut inequality.
The full limiting conclusion above holds with $C=A+a$.
\end{enumerate}
The constants precede the choice of $M$, and $M'$ precedes every choice of
dimension, input sequence, output and time coefficient.
\end{theorem}
\formal{UnconditionalPaidPadding}{TransposeFollowup.PaidPadding.exists_uniform_paid_padding}

\paragraph{What is paid.}
The machine starts with blank work tapes. It derives the required dimension
from the input length, generates its finite schedule, builds the clocks used
to execute it, performs the scheduled excursions, clears the prefix workspace,
resets the input head, and hands off to a faithful simulation of $M$.
All transitions are included in $AP+t$. The implementation uses five work
symbols; the headline contract states their three-bit cell charge. Parallel
bit planes encode the source alphabet using additional tapes, which is why
the displayed tape count depends on the source alphabet while the additive
time coefficient can be chosen before $M$.

The constants are existential; no numerical values are given here. Once the
rational family and its parameters are fixed, the same transition table
handles all positive canonical dimensions.

\begin{proofoutline}
The proof has four parts.

\emph{Analytic existence.} A fixed finite rational family of increasing
partial M\"obius maps expands measurable subsets of the unit interval.
The development follows the continuous-expander strategy of
Bourgain--Yehudayoff~\cite{BourgainYehudayoff2013}: rational free generators
near the identity, flattening and a restricted spectral gap, then the passage
from a gap to expansion. The formal chain uses the pinned scalar-projection
result~\cite{OAIProjection2026} at its additive-combinatorial interface; its
local concentration hypotheses are discharged in our analytic development.
No spectral-gap or expansion hypothesis is left in the wrapper theorem.

\emph{Finite graphs and excursions.} Divide the interval into $d$ cells.
An edge records an overlap of a map image with a target cell. Symmetrization
gives a graph on $2d$ ordered vertices with a fixed positive vertex-expansion
constant. The estimates hold at every positive dimension. Monotonicity gives
a fixed finite page decomposition of these edges. A page is implemented by
one head's excursions, with a constant stretch chosen so that expanding edges
force enough distinct crossing chords in every slot labelling. A fixed number
of extra stationary tapes supplies the required head-count inequality.

\emph{Actual generation and execution.} Rational comparisons determine the
overlap edges and the corresponding finite operation stream. The generator
and its consumer are compiled to finite transition tables. The row-clock and
excursion proofs identify their actual head positions at the actual transition
times, including fetching and setup delays. Their aggregate cost, including
generation, is $O(P)$. Cleanup proves blank tape contents and zero head
positions, and resetting the input head supplies the exact entry contract
for the payload. The five-symbol bit-plane simulation preserves one source
transition per payload transition. Sequential composition therefore preserves
the output and gives the stated additive time bound.

\emph{Extending the horizon and taking limits.} The prefix excursions witness
the cut inequality through the payload's final halting time, using the full
composed machine's head count and paths. Finally
\lean{NativeRunLimits.of_cut} applies the original walk theorem with zero
slack, giving existence of limits and the full assertion for every eligible
limit. The construction parameters are chosen before the payload machine,
as required by the outer quantifiers.
\end{proofoutline}

\subsection{The analytic interface and its attribution}

The following analytic lemma removes the expansion hypothesis from the native
construction. Continuous-expander existence and the proof strategy are due
to Bourgain--Yehudayoff~\cite{BourgainYehudayoff2013}.

\begin{lemma}[Rational continuous-expander interface]\label{lem:rational-expander}
There are a finite nonempty family of rational increasing partial M\"obius
maps $f_j$ and a real $\varepsilon>0$ such that each map has a nonempty open
interval domain with rational endpoints, takes that domain into $[0,1]$, and
\[
 (1+\varepsilon)|S|\le
 \left|\bigcup_j f_j(S\cap\operatorname{dom}f_j)\right|
\]
for every measurable $S\subseteq[0,1]$ with $|S|\le1/2$.
Each map has rational coefficients, positive determinant and positive
denominator at both endpoints of its defining closed interval.
\end{lemma}
\formal{UnconditionalPaidPadding}{TransposeFollowup.exists_rational_continuous_expander}
\formal{RationalMobius}{TransposeFollowup.RationalMobius}

\begin{proofoutline}
The rational free family and a sufficiently fine fixed grid give a restricted
$L^2$ gap for functions supported on $[0,1]$ whose integral on each grid cell
is zero. If a measurable set expanded too little under the family and the
auxiliary translations, its cell-centered, normalized indicator would have
large correlation with every generator and inverse. Averaging contradicts
the restricted gap. Clipping the augmented rational family supplies the
open domains and removes endpoints of measure zero. The same rational
coefficients pass through the gap and expansion interfaces; no separate
unverified approximation of an expander is substituted.
\end{proofoutline}

The lemma and its analytic dependencies are proved in the checked library.
The pinned OpenAI dependency is a formally proved scalar projection result;
it is separate from the OpenAI multiplication source cited in the earlier paper.

\noindent\begin{minipage}{\textwidth}
\paragraph{A short dependency route.}
The following links locate the construction without requiring the reader to
follow the compiler's complete internal file graph.
\begin{center}\small
\begin{tabular}{p{0.31\linewidth}p{0.61\linewidth}}
\toprule
Obligation & Principal interface\\
\midrule
Rational gap and expansion & \lean{UnconditionalPaidPadding}\\
All-size finite expansion & \lean{AllSizeOverlapExpansion}, \lean{ContinuousNativePadding}\\
Generated actual trajectories & \lean{NativeTimedWalk}, \lean{PaidPaddingLimits}\\
Cleanup and blank payload entry & \lean{GeneratorCleanupCost}, \lean{PaidPaddingLimits}\\
Full-horizon cut inequality & \lean{PaidPayload.cut_after_prefix} in \lean{PaidPaddingLimits}\\
Original limiting conclusion & \lean{NativeRunLimits.of_cut}\\
\bottomrule
\end{tabular}
\end{center}
\end{minipage}

\subsection{What the construction establishes}

The earlier paper's expanding-motion proposal~\cite[Construction
programme]{FirstPaper2026} identified preparation and execution as obligations.
Theorem~\ref{thm:padding} resolves those obligations for a wrapper, including
all setup and execution costs. Its construction contribution is a uniform
method for supplying the required motion. It still leaves complete exchange
to the payload: the expander and the excursions do not specify how to recover
the transpose's output bits from coded data.

For the impossibility program, if a linear transposer exists, the theorem
supplies another linear transposer satisfying the zero-slack cut inequality
and the full $M^*$ conclusions. Those necessary conditions therefore cannot
by themselves exclude a linear transposer. An exclusion must also use the
transpose demand or another property not supplied by the wrapper.

The theorem uses $O(P)$ schedule generation. It does not assert an $o(P)$
generator, a particular weak-limit return measure, or a nonhyperfinite return
relation.
